\DocumentMetadata{
  pdfversion=2.0,
  pdfstandard=ua-2,
  lang=en-US,
  tagging=on,
  tagging-setup={math/setup=mathml-SE,tabsorder=structure}
}
\documentclass[11pt,letterpaper]{article}
\usepackage[margin=0.68in,includefoot]{geometry}
\usepackage{newcomputermodern}
\usepackage{amsmath,amscd,mathtools}
\usepackage{tikz,adjustbox}
\providecommand{\square}{\mdlgwhtsquare}
\usepackage[hidelinks]{hyperref}
\hypersetup{
  pdftitle={Algebra First-Year Exam, August 2020, Part 1},
  pdfauthor={University of Florida Department of Mathematics},
  pdfsubject={Graduate mathematics examination},
  pdfdisplaydoctitle=true,
  bookmarksopen=true
}
\setcounter{secnumdepth}{0}
\setlength{\parindent}{0pt}
\setlength{\parskip}{0.28em}
\setlength{\emergencystretch}{3em}
\setlength{\leftmargini}{2.15em}
\setlength{\leftmarginii}{2.65em}
\setlength{\leftmarginiii}{3em}
\pagestyle{empty}
\raggedbottom
\allowdisplaybreaks
\NewTaggingSocketPlug{block/list/label}{examproblem}{%
  \tagstructbegin{tag=Lbl}\tagstructend%
  \tagstructbegin{tag=\LBody}%
  \tagstructbegin{tag=\ProblemHeadingTag,title-o={\ProblemHeadingText}}%
  \tagmcbegin{tag=\ProblemHeadingTag,actualtext={\ProblemHeadingText}}%
  #2%
  \tagmcend%
  \tagstructend%
}
\newenvironment{examproblems}[2]{%
  \def\ProblemHeadingTag{#2}%
  \AssignTaggingSocketPlug{block/list/label}{examproblem}%
  \begin{enumerate}%
  \setcounter{enumi}{#1}%
  \renewcommand{\labelenumi}{\textbf{\theenumi.}}%
  \setlength{\topsep}{0.35em}%
  \setlength{\partopsep}{0pt}%
  \setlength{\itemsep}{0.48em}%
  \setlength{\parsep}{0.18em}%
}{\end{enumerate}}
\newenvironment{examsubparts}{%
  \AssignTaggingSocketPlug{block/list/label}{default}%
  \begin{enumerate}%
  \setlength{\topsep}{0.18em}%
  \setlength{\partopsep}{0pt}%
  \setlength{\itemsep}{0.16em}%
  \setlength{\parsep}{0.1em}%
}{\end{enumerate}}
\newenvironment{examinstructions}{%
  \par\begingroup\itshape
}{%
  \par\endgroup\vspace{0.18em}
}
\makeatletter
\renewcommand\subsection{\@startsection{subsection}{2}{0pt}%
  {-1.25ex plus -.25ex minus -.1ex}%
  {0.55ex plus .1ex}%
  {\normalfont\large\bfseries}}
\renewcommand{\@maketitle}{%
  \newpage\null\vskip -1em%
  {\footnotesize\bfseries
    \href{https://www.ufl.edu/}{University of Florida}, \href{https://math.ufl.edu/}{Department of Mathematics}\par}%
  \vskip -0.5em%
  \tagstructbegin{tag=H1,title={Algebra First-Year Exam, August 2020, Part 1}}%
  \tagpdfparaOff%
  \tagmcbegin{tag=H1}%
    {\LARGE\bfseries\href{https://gma.math.ufl.edu/exams/algebra-first-year/}{Algebra First-Year Exam}, August 2020, Part 1\par}%
  \tagmcend%
  \tagpdfparaOn%
  \tagstructend%
  \vskip 0.25em%
  \tagmcbegin{artifact}%
    \hrule height 1pt%
  \tagmcend%
  \vskip 1em%
}
\makeatother
\title{Algebra First-Year Exam, August 2020, Part 1}
\author{}
\date{}
\begin{document}
\maketitle
\thispagestyle{empty}
\begin{examinstructions}
Answer four problems. Write your answers clearly in complete English sentences.
\end{examinstructions}
\begin{examproblems}{0}{H2}
\def\ProblemHeadingText{Problem 1.}
\item
Denote by \(Q_8\) the quaternion group. Show that if \(f:Q_8\to S_n\) is an injective homomorphism then \(n\geq 8\).
\def\ProblemHeadingText{Problem 2.}
\item
Let~\(G\) be a group.
\begin{examsubparts}
\item[\textnormal{(i)}]
Show that if~\(G\) is abelian, the subset~\(H\) of elements of finite order is a subgroup.
\item[\textnormal{(ii)}]
Give an example of a (nonabelian) group~\(G\) and elements of finite order \(x,y\in G\) such that \(xy\) has infinite order, and prove that this example is correct.
\end{examsubparts}
\def\ProblemHeadingText{Problem 3.}
\item
Prove the first assertion of Sylow’s theorem: if~\(G\) is a group of order~\(n\), \(p\) is prime and \(p^a\) is the largest power of~\(p\) dividing~\(n\), then~\(G\) has a subgroup of order \(p^a\). State clearly any preliminary results you need.
\def\ProblemHeadingText{Problem 4.}
\item
This problem has three parts.
\begin{examsubparts}
\item[\textnormal{(i)}]
Define what it means for a group~\(G\) to be nilpotent.
\item[\textnormal{(ii)}]
Show that any nilpotent group is solvable.
\item[\textnormal{(iii)}]
Give an example of a finite solvable group that is not nilpotent, and prove that this example is correct.
\end{examsubparts}
\def\ProblemHeadingText{Problem 5.}
\item
Let~\(G\) be a group acting on a set~\(S\), and \(H\subseteq G\) a normal subgroup. Assume that for any \(x_1,x_2\in S\) there is a unique \(h\in H\) such that \(h(x_1)=x_2\). For \(x\in S\) let \(G_x\) be the stabilizer of~\(x\). Show that for any \(x\in S\), \(G\) is a semi-direct product of \(G_x\) and~\(H\). Do not assume that~\(G\) is finite.
\def\ProblemHeadingText{Problem 6.}
\item
In both cases, justify your answer.
\begin{examsubparts}
\item[\textnormal{(1)}]
Find all conjugacy classes of elements of order~\(2\) in \(S_6\).
\item[\textnormal{(2)}]
Do the same for \(A_6\).
\end{examsubparts}
\def\ProblemHeadingText{Problem 7.}
\item
Show that the normalizer of a~\(5\)-Sylow subgroup of \(A_5\) has order \(10\), and is a maximal subgroup of \(A_5\).
\end{examproblems}
\end{document}
